% =====================================================================
% Relatorio -- O corte de geracao fotovoltaica no SIN
%
% Um documento so, reunindo as doze analises sobre o corte de geracao
% fotovoltaica no SIN. A forma e a de uma colecao tematica: cada
% capitulo reune as analises que se completam.
%
% ESTRUTURA. Capitulo por bloco de analises que se completam, secao por
% analise. A regra de agrupamento esta escrita no prefacio: analises
% entram no mesmo capitulo quando uma fecha lacuna da outra, e nao por
% compartilharem fonte de dado.
%
% FIGURAS. As doze figuras foram compostas a 183 mm e entram em
% `figure*`, que atravessa as duas colunas: a mancha tem 17,2 cm e a
% reducao e de 6%. Em coluna de 8,25 cm nenhum painel seria legivel.
%
% CLASSE. `book` em duas colunas. O `\part` agrupa: seis partes cabem
% num sumario que o leitor percorre antes de procurar pelo titulo. O
% `\frontmatter` numera o pre-textual em romano, de modo que
% acrescentar pagina ao sumario nao renumera o corpo.
%
% MEDIDAS DA MANCHA: comentario de cabecalho do `estilo-reporte.sty`.
%
% Compilar, da raiz do repositorio:
%   make reporte
% Ou daqui:
%   latexmk -pdf reporte.tex
% O latexmk vive em /Library/TeX/texbin, fora do PATH padrao.
%
% Antes de fechar, rodar os diagnosticos da secao 13 do guia de escrita
% da linha: travessao solto = 0, termos proibidos = 0, mediana de frase
% entre 10 e 30, log sem Overfull nem referencia pendente. O README
% deste diretorio traz os comandos.
% =====================================================================

\documentclass[10pt,a4paper,oneside,twocolumn]{book}

\usepackage[utf8]{inputenc}
\usepackage[brazilian]{babel}
\usepackage[a4paper,top=2.0cm,bottom=2.2cm,left=1.9cm,right=1.9cm,
            columnsep=0.7cm]{geometry}
\usepackage{estilo-reporte}

% A pagina termina onde o texto termina. Com tabela curta em coluna
% estreita, o `\flushbottom` da classe estica a folga entre paragrafos
% ate o rodape e desfaz o ritmo vertical.
\raggedbottom

\hyphenation{sen-so-ria-men-to fo-to-vol-tai-ca fo-to-vol-tai-co
             sub-sis-te-ma in-fle-xi-bi-li-da-de ir-ra-di-an-cia
             de-tec-ta-bi-li-da-de mul-ti-pli-ci-da-de}

\begin{document}

\frontmatter
\onecolumn

% Rosto e sinopse, em uma coluna: a tabela de sinopse tem duas colunas
% de texto e nao cabe em 8,25 cm.
\thispagestyle{empty}

\vspace*{2.2cm}

\begingroup
\sffamily\raggedright
{\fontsize{9}{11}\selectfont\color{tinta2}
Análise de dados abertos do setor elétrico brasileiro \textbullet{} UTFPR\par}
\vspace{1.4em}
{\fontsize{28}{32}\selectfont\bfseries
Relatório\par}
\vspace{0.6em}
{\fontsize{14}{18}\selectfont\color{tinta2}
O corte de geração fotovoltaica no Sistema Interligado Nacional\par}
\vspace{2.2em}
{\fontsize{10}{13}\selectfont
\pe[autor]{nome completo}\par
\vspace{0.3em}
{\color{tinta2}\texttt{opensource.leonardi@gmail.com}}\par
\vspace{0.3em}
{\color{tinta2}setembro de 2026}\par}
\endgroup

\vspace{2.4cm}

\begin{minipage}{15cm}
\setlength{\parskip}{0.6em}
\noindent
Doze análises medem o corte de geração fotovoltaica no Sistema Interligado
Nacional em três eixos: a magnitude do corte, a composição declarada da sua
causa e a sua distribuição entre usinas. Todas usam dados abertos e sem
autenticação, do Operador Nacional do Sistema Elétrico, da ANEEL, do catálogo
público Sentinel-2 e da NASA POWER.

O objetivo inicial era outro. A proposta original detectaria instalações
fotovoltaicas por imagem de satélite para verificar programas de crédito rural.
Duas medidas a redefiniram. A primeira estabelece que 21,6\% da energia
disponível no parque com apuração é cortada, o que separa a geração verificada
do recurso solar. A segunda estabelece que 91\% da micro e minigeração
distribuída ocupa menos de um pixel do Sentinel-2, o que fecha a rota de
verificação por imagem gratuita.

Cinco das doze análises examinam resultados das anteriores. Duas revisam uma
conclusão e três delimitam o escopo de outra. As conclusões revisadas
permanecem no relatório, com a correção correspondente registrada na seção que
a produziu.
\end{minipage}

\vspace{1.6cm}

\begin{minipage}{15cm}
\tabfontw
\begin{tabular}{@{}>{\rag\arraybackslash}p{9.0cm}r@{}}
\toprule
\cab{O que este relatório cobre} & \cab{quantidade} \\
\midrule
Análises com figura e dados-fonte exportados   & 12 \\
Análises que revisam ou delimitam uma anterior &  5 \\
Testes estatísticos que vão a arquivo          & 49 \\
Testes que sobrevivem à correção de Holm       & 16 \\
Conclusões retiradas por análise posterior     &  8 \\
Bases de dados abertas consultadas             &  6 \\
\bottomrule
\end{tabular}
\notas{A correção de multiplicidade é aplicada dentro de cada análise, nunca
entre análises. O apêndice~\ref{ap:inventario} traz janela, amostra, unidade de
repetição e teste declarado de cada uma.}
\end{minipage}

\clearpage

\twocolumn
\tableofcontents

\chapter*{Prefácio}
\addcontentsline{toc}{chapter}{Prefácio}
\markboth{Prefácio}{Prefácio}

\section*{A quem se dirige}

O relatório serve a quem precisa de uma estimativa de exposição ao corte antes
de decidir sobre um ativo fotovoltaico, e a quem modela geração renovável a
partir de irradiância ou de imagem de satélite. Para o primeiro leitor, os
capítulos de despacho e de rede dizem quanto se corta e o que prediz onde se
corta. Para o segundo, o capítulo~\ref{cap:ancora} estabelece que a geração
verificada carrega uma decisão operativa, e não apenas o recurso físico.

Nenhum capítulo pressupõe acesso a dado sob contrato. Todas as fontes são
abertas e sem autenticação, e o apêndice~\ref{ap:procedencia} nomeia cada uma
com o pacote de origem.

\section*{O que fica de fora}

\campo{Micro e minigeração distribuída} O regime de restrição da geração
distribuída é distinto do da geração centralizada, e o corte apurado pelo
operador não a cobre. O capítulo~\ref{cap:satelite} mede a detectabilidade
dessas instalações por satélite, e não o corte delas.

\campo{Ressarcimento} A Lei 15.269/2025 criou compensação por
\textit{constrained-off} durante a execução das análises. O relatório valora a
energia cortada ao Custo Marginal de Operação, que mede valor sistêmico e não
ressarcimento. Os dois divergem por construção.

\campo{Decisão de investimento} O capítulo~\ref{cap:aplicacao} estima quanto
uma área nova produziria e quanto o despacho absorveria. Contrato, outorga e
parecer de acesso não estão em dado aberto e não entram na estimativa.

\campo{Corte anterior a abril de 2024} A apuração pública de
\textit{constrained-off} fotovoltaico começa nessa data. O período anterior é
tratado por fator de capacidade, que responde a várias causas ao mesmo tempo.


\mainmatter

\part{A proposta e o que a redefiniu}
\chapter{A proposta original e as duas medidas que a redefiniram}
\label{cap:ancora}

A proposta que abriu esta linha de trabalho detectaria instalações
fotovoltaicas por imagem de satélite e cruzaria a detecção com o registro da
ANEEL, para verificar programas de crédito rural. Duas medidas a redefiniram
antes que qualquer modelo fosse treinado, e elas fecham caminhos distintos. A
primeira contamina o alvo de qualquer modelo de geração, inclusive de geração
centralizada. A segunda fecha o segmento que a proposta original queria
auditar.

\section{Magnitude e composição do corte no SIN}
\label{sec:fig1}

\begin{ficha}
  \item[Janela]      abril de 2024 a julho de 2026
  \item[Amostra]     85 usinas com apuração de \textit{constrained-off}
  \item[Unidade]     usina por intervalo de 30 minutos
  \item[Fonte]       ONS, restrição por \textit{constrained-off} fotovoltaico
                     e fator de capacidade
  \item[Teste]       nenhum: a análise é descritiva
\end{ficha}

O ONS publica, por intervalo restrito, a geração verificada e uma geração de
referência estimada a partir da curva de desempenho da própria usina. O
\termo{corte} é a diferença entre as duas, truncada em zero:

\[ C = \max\bigl(\text{referência} - \text{verificada},\; 0\bigr) \]

A definição vale apenas nos intervalos em que há razão de restrição apurada. O
campo \texttt{val\_\allowbreak geracaolimitada} não é a energia cortada. Ele registra o teto
ao qual a geração foi limitada. Sua correlação com a geração verificada é 0,86,
e com a perda apurada é 0,02. Usá-lo como energia cortada superestima a perda.

\begin{figure*}[t]
\centering
\includegraphics[width=\linewidth]{../figures/fig1_curtailment_sin.pdf}
\caption{\textbf{O corte de geração fotovoltaica no SIN é sistêmico, de razão
energética e incide sobre toda a janela solar.}
\textbf{a}, Perfil diurno médio da potência fotovoltaica das usinas com apuração
de \textit{constrained-off}. A área azul é a geração verificada e a área
vermelha é a energia cortada; a soma define a potência que estaria disponível.
Médias sobre todos os intervalos de 30 minutos do período, em horário de
Brasília, sem correção para hora solar local.
\textbf{b}, Taxa mensal de corte por subsistema, definida como energia cortada
dividida pela energia disponível.
\textbf{c}, Energia total cortada por razão de restrição, decomposta em origem
sistêmica e local. Os rótulos indicam o total por razão.
\textbf{d}, Distribuição espacial das usinas e conjuntos com apuração. A cor
codifica a taxa de corte acumulada e o tamanho, a capacidade instalada. A escala
de cor satura no percentil 97 para preservar contraste. Contornos estaduais do
IBGE.}
\label{fig:corte-sin}
\end{figure*}

A taxa de corte no parque com apuração é de 21,6\%. A razão energética responde
pela maior parte da energia cortada, e a origem declarada é predominantemente
sistêmica. O corte incide sobre toda a janela solar, e não sobre um horário de
ponta. As taxas do Nordeste e do Sudeste com Centro-Oeste crescem em ritmo
semelhante.

\conclusao{O corte é de razão energética e origem sistêmica, e não de falha
local. Isso o torna uma propriedade do despacho do sistema, e não um atributo
da usina cortada.}

\campo{Consequência para modelagem} Onde a geração de uma usina é limitada por
decisão operativa, um modelo que preveja geração a partir de irradiância ajusta
um alvo que incorpora essa decisão. O mesmo vale para um modelo que parta de
imagem de satélite. O erro não é aleatório, porque a decisão de cortar depende
de carga, de despacho térmico e de posição na rede.

\begin{objecao}
O denominador é o das usinas com apuração, e não o do parque solar nacional. A
geração distribuída tem regime de restrição distinto e não entra. A geração de
referência é estimativa do operador, e em 21,8\% dos intervalos restritos ela
fica abaixo da geração verificada; o truncamento em zero torna a estimativa
conservadora, e não exata. Razão e origem são classificações operativas do ONS,
e não inferência causal independente.
\end{objecao}

\campo{Comparação internacional} Os 21,6\% são de solar centralizada com
apuração. O ONS reporta 20,7\% para eólica e solar em 2025, com denominador
diferente e mesma ordem de grandeza. Os valores publicados para outros sistemas
são menores: 17\% da renovável variável disponível no Chile em 2024, cerca de
11\% na Espanha em julho de 2025, e 4,3\% de média para a solar da Califórnia.
O'Shaughnessy, Cruce e Xu \cite{oshaughnessy2020} organizam a série global de
corte fotovoltaico e mostram que a taxa cresce com a penetração e não com a
idade do sistema, o que põe o valor brasileiro na trajetória esperada e não fora
dela. Novan e Wang \cite{novan2024} medem que a taxa marginal é cerca de duas
vezes a média, o que implica exposição maior para uma usina que entre agora.

\section{Detectabilidade da geração distribuída}
\label{sec:fig2}

\begin{ficha}
  \item[Janela]      conexões registradas até abril de 2026
  \item[Amostra]     4,6 milhões de instalações
  \item[Unidade]     empreendimento
  \item[Fonte]       ANEEL, geração distribuída e informações técnicas
                     fotovoltaicas
  \item[Teste]       nenhum: a análise é descritiva sobre censo administrativo
\end{ficha}

A ANEEL publica a área do arranjo de cada instalação registrada. O número de
pixels que um sensor de resolução $r_s$ entrega sobre um arranjo de área
$A$ é

\[ \mathrm{px}(s) = A / r_s^2 \]

A conta transforma a resolução do sensor em uma grandeza comparável com o
registro, instalação por instalação.

\begin{figure*}[t]
\centering
\includegraphics[width=\linewidth]{../figures/fig2_mmgd_detectabilidade.pdf}
\caption{\textbf{A geração distribuída brasileira é majoritariamente sub-pixel
no Sentinel-2, e a série de conexões carrega uma descontinuidade regulatória.}
\textbf{a}, Distribuição acumulada da área do arranjo das instalações
registradas na ANEEL. As linhas tracejadas marcam a área de um pixel de cada
sensor, e o rótulo indica a fração de instalações menor que um pixel.
\textbf{b}, Fração de instalações que atinge o piso operacional de quatro
pixels, por classe de consumo e sensor.
\textbf{c}, Conexões mensais. A linha tracejada marca o fim da janela de
transição da Lei 14.300/2022, e o marcador indica o mês de pico.
\textbf{d}, Participação da classe rural no total de instalações de cada unidade
da federação. Contornos do IBGE.}
\label{fig:mmgd}
\end{figure*}

No Sentinel-2 a mediana é de 0,29 pixel por instalação, e a fração sub-pixel é
de 91,2\%. No CBERS-4A WPM são 7,25 pixels e 0,5\%. Num sensor comercial de
0,5 m são 116 pixels e fração próxima de zero. A classe rural, alvo dos
programas de crédito, tem área mediana de 41 m\textsuperscript{2} e 5,2\% de
segmentabilidade no Sentinel-2.

\conclusao{Nove em cada dez instalações de geração distribuída ocupam menos de
um pixel do Sentinel-2. A verificação por segmentação exige 2 m ou melhor.}

\campo{A resolução não é a restrição dominante} O registro não traz localização
utilizável. Os campos de coordenada existem no esquema e estão preenchidos em
0,06\% dos casos. Há o campo \texttt{CodCEP}, completo para os 7,7\% de pessoa
jurídica e truncado em cinco dígitos para os 92,3\% de pessoa física. A mediana
é de um prefixo por município, o que não refina o alvo rural. Sem coordenada por
instalação, a verificação exige varredura exaustiva do território. Apenas o
segmento de pessoa jurídica é endereçável.

\campo{A data de conexão não é exógena} A janela de transição da Lei
14.300/2022 concentrou 98 mil conexões em um único mês, quatro vezes a média do
ano anterior. Um estudo de evento sobre a data de conexão mediria a antecipação
regulatória, e não o efeito da política.

\begin{objecao}
Contar pixels não é o mesmo que detectar. Um objeto sub-pixel ainda altera a
assinatura espectral da célula. O que a contagem descarta é a segmentação, e não
a detecção espectral. O piso de quatro pixels é convenção declarada, e não
constante física. O conjunto é censo administrativo, e não amostra, de modo que
não há erro amostral a declarar.
\end{objecao}

\campo{Posição na literatura} Os trabalhos que segmentam telhados operam entre
0,1 e 0,3 m. O DeepSolar \cite{yu2018} usa imagem aérea submétrica, e o conjunto
multirresolução de Jiang et al. \cite{jiang2021} vai de 0,1 a 0,8 m. O
inventário global de Kruitwagen et al. \cite{kruitwagen2021} é a referência com
Sentinel-2, e cobre apenas instalações acima de 10 kW. Não foi localizado
trabalho que segmente telhado residencial a 10 m.

\chapter{Protocolo de teste e regra de multiplicidade}
\label{cap:regra}

Cada análise declara a janela, o tamanho da amostra, a unidade que se repete
dentro dela e o teste aplicado. As declarações completas estão no
apêndice~\ref{ap:inventario}, e cada seção repete a sua na abertura.

\section{Janela, amostra e unidade de repetição}

As janelas de corte saem de uma regra única: percentil 80 do início e percentil
20 do fim do período de apuração de cada fonte. Elas diferem entre fontes porque
a apuração eólica começa em outubro de 2021 e a fotovoltaica em abril de 2024. A
análise da seção~\ref{sec:fig1} vem de extração anterior às demais, e por isso
sua taxa agregada difere em 0,5 ponto da extração final.

A unidade de repetição varia com a pergunta. Nas análises de despacho é a hora
do sistema, com $n$ da ordem de 21 mil. Nas análises de rede é o conjunto de
geração, com $n$ entre 59 e 120. A diferença de duas ordens de grandeza importa
para a leitura: um coeficiente pequeno é detectável no primeiro caso e não é no
segundo.

\section{Correção de multiplicidade}

A correção de Holm é aplicada dentro de cada análise, sobre as oito covariáveis
pré-declaradas dos capítulos de rede, e é declarada ausente onde não foi
aplicada. No estudo de caso pareado as cinco covariáveis testadas resultaram
todas nulas, e a correção só tornaria a conclusão mais conservadora.

Não há correção entre análises. Somados os testes que vão a arquivo, o relatório
reporta 49, dos quais 16 sobrevivem a Holm e 33 não. As oito covariáveis de rede
foram pré-declaradas e testadas na mesma ordem nas duas fontes de geração, e o
capítulo~\ref{cap:rede} submete as sobreviventes a três critérios adicionais. Um
limiar de multiplicidade em nível de relatório não é aplicado.

\begin{table}[t]
\centering
\tabfont
\begin{tabular}{@{}
    >{\rag\arraybackslash}p{2.9cm}
    >{\rag\arraybackslash}p{2.6cm}
    r@{}}
\toprule
\cab{Capítulo} & \cab{Onde entra} & \cab{Testes} \\
\midrule
Estudo de caso        & sem correção      &  5 \\
Detecção por satélite & duas comparações  &  2 \\
Despacho              & quintis de carga  &  5 \\
Decomposição          & sem correção      & 14 \\
Rede, solar           & Holm sobre oito   &  8 \\
Rede, eólica          & Holm sobre oito   &  8 \\
Robustez              & bootstrap, Fisher &  7 \\
\midrule
\cab{Total}           &                   & \cab{49} \\
\bottomrule
\end{tabular}
\caption{Distribuição dos testes que vão a arquivo.}
\label{tab:multiplicidade}
\notas{Dos 49 testes, 16 sobrevivem à correção de Holm aplicada dentro da
análise que os produz. Os 33 restantes são reportados sem correção, e a análise
que os contém declara essa ausência.}
\end{table}

\begin{objecao}
A correção de multiplicidade governa a taxa de erro dentro de cada análise, e
não a validade do resultado. Uma covariável que sobrevive a Holm pode ainda
assim não sobreviver a reamostragem. É o que a seção~\ref{sec:fig12} mede para
os quatro resultados de rede.
\end{objecao}


\part{Observabilidade por satélite}
\chapter{O que o satélite gratuito separa, e a que preço o comercial resolve}
\label{cap:satelite}

A seção~\ref{sec:fig2} fecha a verificação de geração distribuída por imagem
gratuita. Restam duas perguntas: se a rota funciona para geração centralizada,
onde o arranjo tem centenas de hectares, e quanto custaria comprar a resolução
que o sensor gratuito não entrega.

\section{Construção e operação de usina centralizada}
\label{sec:fig7}

\begin{ficha}
  \item[Janela]      2018 a 2026, uma cena por ano
  \item[Amostra]     34 conjuntos, 306 medições
  \item[Unidade]     conjunto por ano
  \item[Fonte]       Sentinel-2 L2A pelo catálogo STAC público; ANEEL/SIGA
  \item[Teste]       Wilcoxon pareado, duas comparações
\end{ficha}

Para cada conjunto mede-se a razão entre o brilho no vermelho dentro do arranjo
e num anel de controle afastado:

\[ \rho(t)=\frac{\overline{B_{\text{red}}}\big|_{d\le 400\,\mathrm{m}}}
{\overline{B_{\text{red}}}\big|_{1{,}5\le d\le 3\,\mathrm{km}}} \]

O anel controla iluminação, estação e umidade do solo, que afetam as duas
regiões da mesma forma. As séries são alinhadas pela data de entrada em operação
registrada no SIGA, e não pelo calendário.

\begin{figure*}[t]
\centering
\includegraphics[width=\linewidth]{../figures/fig7_deteccao_construcao.pdf}
\caption{\textbf{A construção de usinas fotovoltaicas é detectável e datável por
Sentinel-2; a operação não é.}
\textbf{a}, Razão entre o brilho no vermelho dentro do arranjo e num anel de
controle de 1,5 a 3 km, por usina em linha cinza e mediana entre usinas em linha
vermelha, com faixa entre quartis, alinhadas pela data de entrada em operação do
SIGA.
\textbf{b}, Sentinel-2 em cor natural sobre o complexo Sol do Cerrado, em escala
de refletância fixa entre as datas. Os círculos marcam as 17 usinas do SIGA.
\textbf{c}, Comparação pareada da razão de brilho entre fases, uma linha por
usina, com a mediana em destaque, e valor de $p$ de Wilcoxon pareado.
\textbf{d}, Diferença entre o ano de pico do brilho relativo e o ano de entrada
em operação. As barras em vermelho indicam erro de até um ano.}
\label{fig:construcao}
\end{figure*}

Da fase de base para a de construção a razão sobe de 1,01 para 1,30, um ganho de
34\% na mediana, em 24 de 30 usinas, com $p = 1{,}2 \times 10^{-5}$. Da base
para a operação a mediana vai de 0,95 para 1,02, um ganho de 6,6\%, com 11 de 18
usinas em alta contra 7 em queda e $p = 0{,}09$.

\conclusao{A obra tem assinatura espectral e data a entrada em operação. A fase
operacional não se separa do entorno a 10 m de resolução.}

A construção remove vegetação de centenas de hectares, e solo exposto tem
refletância alta no vermelho. Na fase operacional não há sinal em nenhuma
direção. Módulos são escuros, mas a 10 m o pixel mistura painel, corredor, via
de serviço e solo entre fileiras.

\campo{O que isso responde da proposta original} A hipótese de cruzar registro
com evidência de satélite é executável para geração centralizada, com uma
inversão. O satélite confirma a usina pela cicatriz da obra, e não pelo arranjo
em si. Como a cicatriz é evento transitório, a auditoria exige monitoramento
contínuo da série, e não observação do estado atual.

\campo{Uma correção de método dominou o resultado} A primeira execução usou a
coordenada do conjunto publicada pelo ONS, que corresponde à subestação
coletora, e resultou nula. Em 31 de 37 conjuntos a subestação fica a mais de
750 m do centroide do arranjo, com mediana de 1,9 km. A máscara amostrava
terreno fora do arranjo. Refeita com as coordenadas de usina do SIGA, a
diferença entre fases aparece.

\begin{objecao}
A qualidade da geolocalização de referência domina o resultado, e qualquer
replicação depende dessa camada. Uma cena por ano fixa a datação em resolução
anual, com acerto exato em 21\% dos casos e erro de até um ano em 62\%. Os
círculos de 400 m não são o polígono do arranjo, e delinear por segmentação
melhoraria o contraste. O pulso indica remoção de vegetação, que também ocorre
em lavoura, mineração e loteamento; a atribuição depende do cruzamento com o
registro. Nada disso se estende a telhados.
\end{objecao}

\campo{Posição na literatura} A razão dentro e fora do arranjo é uma versão
simplificada do \textit{Temporal Cluster Matching} \cite{robinson2021}, que data
construções comparando a relação espectral dentro e fora da pegada. Aquele
método foi aplicado a usinas solares em Karnataka. O Global Renewables Watch
\cite{grw2025} estima datas de construção em escala global com imagem de alta
resolução. O que esta seção acrescenta é a validação contra o registro de
outorga e o resultado negativo sobre a fase operacional.

\section{Resolução, cobertura e custo de imagem comercial}
\label{sec:aux}

A curva de detectabilidade da seção~\ref{sec:fig2} posiciona qualquer catálogo
de sensor comercial. Os 21 sensores de um fornecedor típico caem sobre ela em
três grupos.

\begin{figure*}[t]
\centering
\includegraphics[width=\linewidth]{../figures/aux_geopera_capacidade.pdf}
\caption{\textbf{Os sensores de um catálogo comercial posicionados contra a
curva de detectabilidade da geração distribuída, e o custo de aquisição por
área.}
\textbf{a}, Resolução e faixa de imageamento dos sensores do catálogo.
\textbf{b}, Fração da geração distribuída registrada que atinge o piso de quatro
pixels em cada resolução, com os sensores do catálogo sobrepostos.
\textbf{c}, Custo de aquisição por instalação verificada, no cenário com
coordenada por instalação e no cenário de varredura municipal.}
\label{fig:comercial}
\end{figure*}

A transição de detectabilidade ocorre entre 1 m e 5 m. Abaixo de 1 m, 99,5\% da
geração distribuída é segmentável, e os 18 sensores submétricos do catálogo
estão todos no platô. Acima, a fração cai para 86,6\% a 2 m, 7,5\% a 5,3 m e
1,6\% no Sentinel-2. Os dois sensores hiperespectrais do catálogo ficam acima do
joelho da curva.

A faixa de 2 m custa entre AUD 2 e 3 por km\textsuperscript{2} e entrega a mesma
resolução do CBERS-4A WPM, que é público. Entre 0,80 m e 2 m há apenas arquivo,
sem programação de aquisição.

\conclusao{O custo de verificação é determinado pela ausência de localização no
registro, e não pela resolução disponível no mercado.}

Com coordenada por instalação, um recorte de 0,01 km\textsuperscript{2} a 30 cm
custa entre AUD 0,18 e 0,28. Sem ela, a varredura de 1.000 km\textsuperscript{2}
de um município custa entre AUD 18 mil e 28 mil. São cinco ordens de grandeza
entre os dois cenários.

\campo{O que fica viável} Varrer entre 100 e 400 km\textsuperscript{2} da área
ocupada de dois ou três municípios rurais custa entre AUD 1.800 e 11.200, e
permite auditar a contagem agregada contra o registro municipal. A verificação
por instalação na classe rural permanece inviável.

As 361.622 instalações de pessoa jurídica têm CEP completo no registro,
geocodificável em nível de rua. Para esse segmento, com arranjos medianos de 66
e 92 m\textsuperscript{2} em comércio e indústria, o recorte dirigido é viável,
com custo da ordem de AUD 0,20 por instalação.

\begin{objecao}
Preços e sensores são os publicados pelo fornecedor, em dólar australiano e sem
tributo, e mudam sem aviso. A profundidade de arquivo não é publicada por
sensor, o que é decisivo porque todo desenho deste relatório é longitudinal. O
catálogo é apenas óptico, sem radar de abertura sintética, que seria o
complemento imune a nuvem.
\end{objecao}

\campo{Onde está a análise energética por satélite} A abordagem em uso
operacional emprega satélite meteorológico geoestacionário, com índice de nuvem
convertido em irradiância e depois em geração por modelo de desempenho, a 1 ou
2 km e passo de 10 minutos. Ela mede a nuvem, e não o arranjo. É a abordagem do
capítulo~\ref{cap:calibracao}.


\part{Restrição ou recurso}
\chapter{Estudo de caso pareado e cronologia do entorno}
\label{cap:caso}

Duas usinas do mesmo subsistema, sujeitas ao mesmo despacho, perdem frações
distintas da energia disponível. As explicações locais que a diferença sugere
são testáveis na frota, e o período anterior ao início da apuração é alcançável
pelo fator de capacidade.

\section{Dois ativos com taxas de corte distintas}
\label{sec:fig3}

\begin{ficha}
  \item[Janela]      outubro de 2024 a agosto de 2026, recorte comum ao par
  \item[Amostra]     2 ativos no par; 60 conjuntos no teste de frota
  \item[Unidade]     conjunto
  \item[Fonte]       ONS e ANEEL/SIGA
  \item[Teste]       Spearman sobre 5 covariáveis, sem correção
\end{ficha}

Sol do Cerrado tem 681,3 MW, é produtor independente, reúne 17 sociedades de
propósito específico de um mesmo grupo e conecta em barramento compartilhado de
Jaíba, em 230 kV. O conjunto Três Marias 3 tem 70,0 MW, opera sob autoprodução
de sete empresas privadas distintas e tem ponto de conexão exclusivo. Os dois
ficam em Minas Gerais, no mesmo subsistema. A janela comum foi recortada
programaticamente.

\begin{figure*}[t]
\centering
\includegraphics[width=\linewidth]{../figures/fig3_estudo_caso.pdf}
\caption{\textbf{Dois ativos solares em Minas Gerais diferem 3,0 vezes na taxa
de corte, e a diferença não se reproduz como padrão no parque.}
\textbf{a}, Perfil diurno médio de potência, em pu da capacidade instalada, na
janela comum aos dois ativos. Em azul a geração verificada, em vermelho a
energia cortada.
\textbf{b}, Custo Marginal de Operação do subsistema Sudeste ponderado pela
energia, durante os intervalos de corte e durante a geração.
\textbf{c}, Taxa de corte mensal dos dois ativos.
\textbf{d}, Pixels Sentinel-2 ocupados pela área de módulos, comparados à
instalação mediana de geração distribuída, em escala logarítmica, com área de
módulos estimada a 5,12 m\textsuperscript{2}/kWp.
\textbf{e}, Taxa de corte contra capacidade renovável variável num raio de
100 km, para os conjuntos com apuração completa na janela de referência da
frota, com os dois alvos em destaque.}
\label{fig:caso}
\end{figure*}

A taxa de corte é de 26,8\% contra 8,8\%. A energia cortada equivale a cerca de
R\$ 64 milhões contra pouco mais de R\$ 1 milhão, valorada ao Custo Marginal de
Operação. Nos dois ativos o CMO ponderado durante o corte é menor que durante a
geração, o que é consistente com restrição de razão energética. O
capítulo~\ref{cap:despacho} examina esse ponto no despacho.

\campo{Uma correção de atribuição} O campo \texttt{nom\_\allowbreak agenteproprietario} do
ONS identifica o agente representante, e não o proprietário. Ele não sustenta a
leitura de que Três Marias 3 seja ativo estatal. O SIGA registra sete
proprietários privados sob regime de autoprodução. Não há usina solar estatal no
parque despachado, e o eixo público contra privado não é testável nesta amostra.

\conclusao{Nenhuma das cinco covariáveis locais reproduz a diferença do par na
frota: nem carga do ponto de conexão, nem número de conjuntos, nem densidade
regional de renovável, nem tensão, nem capacidade própria.}

Duas explicações locais eram plausíveis à vista do par: barramento compartilhado
e densidade regional de renovável. Nenhuma sobrevive ao teste na frota. Uma
explicação consistente com um par de ativos não se generaliza a mecanismo sem
esse teste.

\begin{objecao}
Dois ativos não são amostra, e o par foi escolhido por contraste. As cinco
covariáveis foram testadas sem correção de multiplicidade; como nenhuma
resultou significativa, a correção só tornaria a conclusão mais conservadora. O
CMO é proxy de valor sistêmico, e não preço de liquidação nem prejuízo líquido.
\end{objecao}

\campo{O que estava em dado aberto e foi afirmado em falta} A versão anterior
desta análise concluía que a topologia de transmissão não está publicada. Ela
está. O SIGET publica o grafo linha a linha, com 1.166 linhas e 619 subestações
com código ONS, tensão e comprimento, e o pacote \texttt{fator-\allowbreak capacidade-2}
traz a coordenada do ponto de conexão de cada usina. O join cobre 168 dos 229
conjuntos despachados. O limite de transferência também está publicado, no
pacote \texttt{linha-\allowbreak transmissao}, com capacidade operativa em MVA e variantes
sazonais, além de resistência e reatância. Permanece indisponível apenas qual
restrição está ativa em cada intervalo. As cinco covariáveis testadas são locais
ao ponto e não usam a estrutura da rede, de modo que o resultado nulo não se
estende a ela. O capítulo~\ref{cap:rede} a testa.

\section{Cronologia do corte e da energização do entorno}
\label{sec:fig4}

\begin{ficha}
  \item[Janela]      2023 a julho de 2026; entorno desde 2021
  \item[Amostra]     43 e 24 meses, por ativo
  \item[Unidade]     conjunto por mês
  \item[Fonte]       ONS, fator de capacidade e \textit{constrained-off}
  \item[Teste]       nenhum teste formal
\end{ficha}

A apuração de corte começa em abril de 2024, e a onda de construção em Jaíba vai
de 2021 a 2024. A série de corte não cobre o período em que a hipótese de
saturação do entorno se manifestaria. O fator de capacidade existe desde a
entrada em operação e cai quando a geração é limitada:

\[ FC_{\text{verificado}}(t)=\frac{E_{\text{gerada}}(t)}{P_{\text{inst}}\cdot\Delta t} \]

A partir de abril de 2024 a apuração permite recompor o potencial:

\[ FC_{\text{potencial}} = FC_{\text{verificado}} + FC_{\text{cortado}} \]

Se a queda for de restrição de rede, o potencial recomposto permanece estável
enquanto o verificado cai. Se for recurso solar ou degradação, os dois caem
juntos.

\begin{figure*}[t]
\centering
\includegraphics[width=\linewidth]{../figures/fig4_cronologia.pdf}
\caption{\textbf{A queda do fator de capacidade de Sol do Cerrado acompanha a
energização do entorno; o controle não mostra o mesmo.}
\textbf{a}, Jaíba, Minas Gerais. Em área cinza a capacidade fotovoltaica
acumulada em operação num raio de 100 km, no eixo esquerdo; em linha o fator de
capacidade mensal verificado de Sol do Cerrado, no eixo direito; em faixa
vermelha a parcela cortada apurada pelo ONS, disponível a partir de abril de
2024.
\textbf{b}, São Gonçalo do Abaeté, Minas Gerais, com a mesma construção para
Três Marias 3.
\textbf{c}, Fator de capacidade mensal de Sol do Cerrado em 2023 e 2025, mês
contra mês, para separar restrição de sazonalidade.
\textbf{d}, Potência acumulada de geração distribuída nos dois municípios. Os
20,4 MW de Jaíba equivalem a 3,0\% da capacidade de Sol do Cerrado.}
\label{fig:cronologia}
\end{figure*}

\begin{objecao}
A decomposição separa restrição de degradação, e não separa restrição de
variação climática. Se 2025 tiver sido menos ensolarado que 2023, a queda não
informa sobre a rede. O fator de capacidade responde a várias causas ao mesmo
tempo, e a decomposição só vale onde a apuração existe. A comparação é pareada e
sem aleatorização, como na seção~\ref{sec:fig3}.
\end{objecao}

O capítulo~\ref{cap:calibracao} fecha essa lacuna, medindo a irradiância
diretamente sobre cada usina.

\campo{Por que o desenho não estima efeito de tratamento} Cada vizinho energiza
em data distinta, o que caracteriza adoção escalonada. Callaway e Sant'Anna
\cite{callaway2021} mostram que o estimador de diferenças em diferenças com
efeitos fixos duplos fica enviesado nesse caso, por ponderar comparações entre
unidades já tratadas. A leitura é descritiva e pareada contra um único
controle, o que evita o problema e também não estima efeito. Estimá-lo exigiria
o estimador por grupo e período daqueles autores, e uma definição de tratamento
exógena que a seção~\ref{sec:fig2} mostra não existir do lado da geração
distribuída.

\chapter{Controle climático e calibração da referência}
\label{cap:calibracao}

O capítulo anterior mede uma queda de fator de capacidade sem separar restrição
de variação do recurso. A irradiância medida sobre cada usina faz essa
separação, e a separação depende de um pressuposto sobre o ano de base.

\section{Restrição ou recurso}
\label{sec:fig5}

\begin{ficha}
  \item[Janela]      2023 a julho de 2026
  \item[Amostra]     meses com cobertura climática
  \item[Unidade]     conjunto por mês
  \item[Fonte]       ONS; NASA POWER, irradiância diária no ponto da usina
  \item[Teste]       nenhum: decomposição aritmética
\end{ficha}

Com a irradiância global horizontal diária $H$ sobre a usina, a razão de
desempenho divide a energia gerada por kW instalado pelo recurso disponível:

\[ PR = \frac{FC \cdot 24\,\mathrm{h}}{H} \]

Se a queda do fator de capacidade viesse de menor irradiância, $PR$ permaneceria
estável.

\begin{figure*}[t]
\centering
\includegraphics[width=\linewidth]{../figures/fig5_controle_climatico.pdf}
\caption{\textbf{Normalizada pela irradiância medida sobre a usina, a perda de
Sol do Cerrado permanece; o recurso solar não a explica.}
\textbf{a}, Razão de desempenho mensal, definida como energia gerada por kW
instalado dividida pela irradiância global horizontal diária. O tracejado
horizontal marca a média de cada ano.
\textbf{b}, Irradiância global horizontal mensal sobre cada usina, em linha
cheia, e a mesma sob céu claro, em tracejado.
\textbf{c}, Decomposição da queda do fator de capacidade de Sol do Cerrado entre
2023 e 2025: observado em 2023, esperado em 2025 caso apenas a irradiância
tivesse mudado, e observado em 2025.
\textbf{d}, Fator de capacidade mensal contra irradiância mensal em Sol do
Cerrado, por ano.}
\label{fig:clima}
\end{figure*}

A irradiância sobre Jaíba caiu 2,8\% entre 2023 e 2025, e o fator de capacidade
caiu 29,1\%. Escalando o fator de 2023 pela razão de irradiância, o esperado
para 2025 seria 0,206, e o observado foi 0,150. Cerca de um décimo da queda é
atribuível ao recurso. A razão de desempenho cai de 0,854 para 0,619 em Sol do
Cerrado, enquanto o controle Três Marias 3 permanece em torno de 0,99 sem
tendência.

\conclusao{Nove décimos da queda do fator de capacidade não são atribuíveis à
irradiância. A perda é de rede, e não de recurso.}

Os meses de 2025 não estão sobre a curva de 2023 em posição de menor
irradiância. Eles estão sobre uma curva mais baixa. O deslocamento é consistente
com perda independente do recurso.

\campo{Uma conclusão secundária que não se sustentou} O déficit não explicado
pela irradiância é de 0,056 em fator de capacidade, e o corte apurado pelo ONS
equivale a 0,147, ou 2,6 vezes mais. A diferença foi interpretada como indício
de que a geração de referência do operador superestima o contrafactual. A
interpretação pressupõe que 2023 fosse um ano sem restrição. A
seção~\ref{sec:fig6} mostra que não era.

\begin{objecao}
A irradiância vem de reanálise, e não de piranômetro: NASA POWER combina MERRA-2
com satélite em célula de 0,5 por 0,625 grau. A grandeza é irradiância
horizontal, e os módulos não são horizontais. A razão de desempenho é indicador
relativo dentro de cada usina, e não a razão de norma IEC. A
temperatura não é modelada e o contrafactual é linear.
\end{objecao}

\campo{Erro de variável derivada de sensoriamento} A irradiância entra como
denominador. Proctor, Carleton e Sum \cite{proctor2023} mostram que uma variável
assim enviesa o ponto e o erro padrão de uma inferência a jusante, mesmo com
erro de medida pequeno. Eles recomendam tratá-la como distribuição, por
imputação múltipla. Alix-García e Millimet \cite{alixgarcia2023} acrescentam que o erro de
produto de satélite raramente é clássico. Nenhuma das duas correções foi
aplicada. A diferença medida é grande o bastante para sobreviver a erro
moderado, e o intervalo em torno dela é mais estreito do que deveria ser.

\section{Calibração da geração de referência}
\label{sec:fig6}

\begin{ficha}
  \item[Janela]      abril de 2024 a julho de 2026
  \item[Amostra]     74 usinas com 30 dias limpos ou mais
  \item[Unidade]     usina; dia-usina na calibração
  \item[Fonte]       ONS; NASA POWER
  \item[Teste]       Spearman sobre 2 covariáveis
\end{ficha}

Para cada usina calibra-se uma razão de desempenho nos dias sem restrição da
própria usina, com fração cortada abaixo de 1\%. Compara-se então o potencial
implicado pelo ONS, igual à geração verificada mais a cortada, com o potencial
previsto pela irradiância local sob aquela razão calibrada.

\begin{figure*}[t]
\centering
\includegraphics[width=\linewidth]{../figures/fig6_vies_referencia.pdf}
\caption{\textbf{A geração de referência do ONS é bem calibrada; o viés aparente
da seção anterior vinha do contrafactual, e não do operador.}
\textbf{a}, Distribuição do viés mediano por usina, definido como a razão entre
o potencial implicado pelo ONS e o potencial previsto pela irradiância local sob
a razão de desempenho calibrada nos dias sem restrição da própria usina. O
tracejado em 1 marca ausência de viés.
\textbf{b}, Viés contra a intensidade mediana de corte nos dias restritos.
\textbf{c}, Validação da calibração: fator de capacidade contra irradiância nos
dias sem restrição de Sol do Cerrado, com a reta ajustada pela razão de
desempenho limpa.
\textbf{d}, Razão de desempenho anual de Sol do Cerrado contra o desempenho
limpo calibrado. O rótulo interno indica o corte implícito de cada ano.}
\label{fig:calibracao}
\end{figure*}

O viés mediano por usina é de 1,02, com metade das usinas entre 0,98 e 1,06.
Nenhuma das 74 usinas calibradas passa de 1,5, e 26 ficam abaixo de 1, o que
torna a referência conservadora para um terço delas. No agregado energético dos
dias restritos, o corte apurado é 1,03 vezes o implicado pela irradiância. A
associação entre viés e intensidade de corte é nula, o que descarta a falha de
maior consequência.

\conclusao{A geração de referência do ONS não apresenta viés detectável, e não
se degrada onde o corte é mais intenso.}

\campo{A revisão} O desempenho limpo de Sol do Cerrado, calibrado em 279 dias
sem restrição, é de 0,931. Em 2023 a usina operou a 0,854, ou 8,3\% abaixo do
seu próprio desempenho limpo. Aquele ano já era restrito, e não havia apuração
pública. Ao usar 2023 como base, a seção~\ref{sec:fig5} mediu um déficit
reduzido e concluiu de forma incorreta que o operador superestimava o corte.

\campo{O que muda e o que permanece} A conclusão principal da
seção~\ref{sec:fig5} permanece e sai reforçada. Contra o desempenho limpo, o
corte implícito é de 33\% em 2025, e não de 26\%. A afirmação secundária sobre
viés do operador é retirada. A taxa agregada da seção~\ref{sec:fig1} não requer
correção, porque uma referência calibrada sustenta o agregado.

\begin{objecao}
Dias classificados como limpos podem não ser limpos. Como o resultado é ausência
de viés, esse erro só o tornaria mais forte. Seis usinas ficaram de fora por não
atingirem 30 dias limpos, e são as mais restritas, onde a referência é mais
usada. A razão de desempenho limpa é constante por usina e herda qualquer viés
sazonal dos dias de calibração.
\end{objecao}

\campo{Posição na literatura} A geração de referência do ONS é derivada de
curvas de desempenho da própria usina, como descrevem Vieira et al.
\cite{vieira2026}. Validá-la exige fonte independente. No lado eólico, Silva et
al. \cite{silva2026} usam anemometria de SCADA. A verificação equivalente no
lado solar, com irradiância de satélite, não foi localizada na literatura
consultada, e a busca não foi sistemática.


\part{Despacho e alocação do corte}
\chapter{Despacho fora da ordem de mérito}
\label{cap:despacho}

A seção~\ref{sec:fig3} registra que o Custo Marginal de Operação durante o corte
é menor que durante a geração. O sistema corta quando a energia vale pouco. A
decomposição de despacho publicada pelo operador diz com o que esse corte
coincide no mesmo instante.

\section{Coincidência entre corte e térmica inflexível}
\label{sec:fig8}

\begin{ficha}
  \item[Janela]      abril de 2024 a agosto de 2026
  \item[Amostra]     21.192 horas
  \item[Unidade]     hora do SIN
  \item[Fonte]       ONS: \textit{constrained-off}, geração térmica por motivo
                     de despacho, curva de carga, energia armazenada
  \item[Teste]       Mann-Whitney e Spearman, estratificado por quintil de carga
\end{ficha}

A decomposição de geração térmica do ONS é aninhada, e não aditiva. Somar todos
os campos dá 1,58 vez a geração total. A definição consistente subtrai a parcela
econômica do total:

\begin{equation*}
\begin{aligned}
\text{fora do mérito} = {} & \texttt{verifgeracao} \\
                          {}-{} & \texttt{verifordemmerito}
\end{aligned}
\end{equation*}

A definição foi verificada contra a soma dos motivos não econômicos disjuntos,
com razão de 1,0030 e mediana da diferença de 0,17 MWmed.

\begin{figure*}[t]
\centering
\includegraphics[width=\linewidth]{../figures/fig8_despacho_merito.pdf}
\caption{\textbf{O corte fotovoltaico por excesso de oferta coincide com térmica
inflexível, mas apenas em carga baixa.}
\textbf{a}, Perfil diurno médio nacional. Em área azul a geração fotovoltaica
verificada, em área vermelha empilhada a energia cortada por razão energética, e
em linhas a geração térmica separada em parcela econômica, por ordem de mérito,
e parcela fora do mérito. A anotação registra a carga líquida, igual à carga do
SIN menos a geração fotovoltaica, na madrugada e no pico solar.
\textbf{b}, Composição da geração térmica fora do mérito nas horas com e sem
corte por razão energética, decomposta pelos motivos não econômicos disjuntos do
dicionário do ONS. O rótulo indica o total.
\textbf{c}, Correlação de Spearman entre corte por razão energética e geração
térmica fora do mérito, estratificada por quintil de carga do SIN, restrita às
horas de 9 h a 15 h. As barras cinza indicam associação não significativa, e
abaixo do eixo aparece o corte médio de cada quintil.
\textbf{d}, Controle negativo: taxa diária de corte por razão energética contra
a energia armazenada do SIN, com a mediana por quartil de armazenamento.}
\label{fig:despacho}
\end{figure*}

A janela acumula 21.889 GWh de corte, dos quais 68,5\% por razão energética. Nas
horas com esse corte há 4.564 MWmed de térmica fora do mérito, contra 3.146 nas
demais, um acréscimo de 45\%. Ponderado pela energia cortada, o valor sobe para
5.266 MWmed, dos quais 73\% são inflexibilidade pura.

Da madrugada ao pico solar a parcela por mérito cai 1.640 MWmed e a parcela fora
do mérito sobe 1.184 MWmed. A mesma térmica permanece ligada e deixa de ser
econômica quando o CMO cai.

\conclusao{Dois terços do corte são classificados como excesso de oferta e
ocorrem com 4,6 GWmed de térmica fora da ordem de mérito no mesmo instante.}

\campo{A leitura que o controle de carga desfaz} A hipótese de que a térmica
desloca a geração solar não sobrevive ao controle de carga. A carga sobe
aproximadamente o quanto a geração solar sobe, e a carga líquida é plana, de
72,2 para 72,1 GWmed. O que varia ao longo do dia é a classificação do despacho
térmico, e não o seu volume.

\campo{A associação é condicional} A correlação entre corte e térmica inflexível
é de 0,48 no quintil de carga mais baixa e de 0,38 no segundo, e cai para a
faixa de 0,09 a 0,20 nos três superiores. Com $n$ da ordem de 1.240 por estrato,
todos os quintis dão $p < 0{,}05$. A leitura se apoia na magnitude, e não na
significância, e a queda não é monotônica.

\begin{objecao}
Coincidência não é causalidade. Separar as duas exigiria reotimizar o despacho
sem a inflexibilidade declarada. A inflexibilidade é declaração do agente, e não
mínimo técnico verificável em dado aberto. O agregado é nacional, e a
seção~\ref{sec:fig9} o revisa. Como controle negativo, a energia armazenada não
explica o corte, com $\rho = 0{,}06$ e $p = 0{,}06$.
\end{objecao}

\campo{Posição na literatura} O mecanismo é documentado na literatura chinesa de
\textit{curtailment}. Unidades térmicas com mínimo técnico e status
\textit{must-run} ocupam a faixa que a renovável variável utilizaria, e a
resposta de política foi o programa de \textit{flexibility retrofit}
\cite{cem2018}, com 16,35 GW em demonstração em 2017. No Brasil, Vieira et al.
\cite{vieira2026} modelaram o corte com regressão logística horária. Eles
identificaram a geração distribuída como preditor mais forte do corte por
demanda, com razão de chances próxima de 9,8, e a carga do sistema atuando de
forma protetiva. É o gradiente que o painel \textbf{c} reproduz.

\section{Decomposição por subsistema}
\label{sec:fig9}

\begin{ficha}
  \item[Janela]      a mesma, restrita às horas de 9 h a 15 h
  \item[Amostra]     21.192 horas
  \item[Unidade]     hora do subsistema
  \item[Fonte]       ONS, com intercâmbio nacional acrescentado
  \item[Teste]       Spearman, sem correção de multiplicidade
\end{ficha}

Nordeste e Sudeste com Centro-Oeste cortam volumes semelhantes nas horas
solares. A térmica inflexível não acompanha.

\begin{table}[t]
\centering
\tabfont
\begin{tabular}{@{}>{\rag\arraybackslash}p{3.6cm}rr@{}}
\toprule
\cab{MWmed, 9 h a 15 h} & \cab{NE} & \cab{SE/CO} \\
\midrule
corte por razão energética & 1.079 & 950 \\
inflexibilidade pura local & 36 & 1.739 \\
razão corte sobre piso local & 30,2 & 0,5 \\
$\rho$ com térmica do NE & +0,08 & +0,10 \\
$\rho$ com térmica do SE & +0,42 & +0,41 \\
\bottomrule
\end{tabular}
\caption{Corte e piso térmico inflexível por subsistema.}
\label{tab:subsistema}
\notas{A razão corte sobre piso local mede quantas vezes o corte do subsistema
excede a inflexibilidade pura declarada nele. As duas últimas linhas medem a
mesma correlação contra a térmica de cada subsistema.}
\end{table}

\begin{figure*}[t]
\centering
\includegraphics[width=\linewidth]{../figures/fig9_decomposicao_subsistema.pdf}
\caption{\textbf{O piso térmico inflexível está onde o corte não está, e a
coincidência nacional não sustenta leitura causal local.}
\textbf{a}, Potência média nas horas solares por subsistema: corte fotovoltaico
por razão energética, geração térmica fora da ordem de mérito e a parcela desta
que é inflexibilidade pura. Abaixo do eixo, a razão entre o corte e o piso
inflexível local.
\textbf{b}, Correlação de Spearman entre corte por razão energética e térmica
fora do mérito do mesmo subsistema, por quintil de carga do próprio subsistema,
com tracejado no limiar de magnitude declarado.
\textbf{c}, A mesma correlação medida contra a térmica do próprio subsistema e
contra a do outro subsistema com corte material.
\textbf{d}, Distribuição acumulada do fluxo Nordeste para Sudeste nas horas de
corte alto e nas demais, com tracejado no percentil 95 do fluxo, usado como
referência empírica de teto por não haver limite operativo publicado.}
\label{fig:subsistema}
\end{figure*}

No Nordeste, onde está o maior volume de corte do país, há 36 MWmed de
inflexibilidade contra 1.079 de corte. A associação local é de 0,08 e não passa
de 0,16 em nenhum quintil de carga.

No Sudeste com Centro-Oeste o painel \textbf{b} reproduz o padrão da
seção~\ref{sec:fig8}, com $\rho = 0{,}40$ nos quintis de carga baixa. O painel
\textbf{c} não sustenta a leitura local. A térmica do Sudeste se correlaciona
com o corte do Sudeste em 0,41 e com o do Nordeste em 0,42, a mais de dois mil
quilômetros de distância. O valor de $\rho$ é determinado por qual térmica se
mede, e não por qual corte. Uma causa local produziria associação mais forte com
o corte local.

\conclusao{A coincidência nacional entre corte e térmica inflexível é real e
compatível com um fator comum, a carga nacional baixa. Este desenho não distingue
as duas hipóteses, e a atribuição causal da seção anterior é retirada.}

\campo{A hipótese de corredor também não se sustenta} Nas horas de corte alto o
fluxo do Nordeste para o Sudeste é menor, de 2.993 contra 3.848 MWmed. Ele fica
acima do teto empírico em 2,9\% das horas, contra 5,7\% nas demais, com
$\rho = -0{,}25$. O Nordeste é cortado com o corredor operando abaixo do seu
percentil habitual.

\campo{O que muda na seção anterior} Nenhum valor da seção~\ref{sec:fig8} muda.
O que se retira é a atribuição causal.

\begin{objecao}
Subsistema ainda é agregação grosseira, e o Sudeste com Centro-Oeste vai de
Minas ao Mato Grosso do Sul. O teto do corredor é o percentil 95 do próprio
fluxo. Não há limite operativo publicado: o pacote
\texttt{programacao\_\allowbreak fluxo\_\allowbreak controlado} cobre apenas
elos de corrente contínua e interligações internacionais, e o pacote
\texttt{linha-\allowbreak transmissao} dá capacidade por linha, e não do
corredor, que resulta de estudo de estabilidade. A ausência de demanda no destino é inferência
por eliminação, e não foi testada de forma positiva. O painel \textbf{c} não
estabelece ausência de mecanismo local no Sudeste; ele mostra que este desenho
não o identifica, porque inflexibilidade local e carga nacional são colineares.
\end{objecao}

\chapter{Topologia da rede e alocação do corte entre usinas}
\label{cap:rede}

A seção~\ref{sec:fig3} testou cinco covariáveis locais ao ponto de conexão e não
achou nenhuma. Restam as covariáveis estruturais, que só existem se houver um
grafo de transmissão montado: distância até a carga, grau do nó, número de
usinas no ponto e outras cinco.

\section{Oito covariáveis de rede na geração fotovoltaica}
\label{sec:fig10}

\begin{ficha}
  \item[Janela]      setembro de 2025 a agosto de 2026
  \item[Amostra]     59 conjuntos fotovoltaicos
  \item[Unidade]     conjunto
  \item[Fonte]       EPE, linhas e subestações em operação; ONS
  \item[Teste]       Spearman sobre 8 covariáveis, com correção de Holm
\end{ficha}

A rede da EPE traz 920 subestações e 1.896 linhas resolvidas em aresta, ou 87\%
do total, com ancoragem mediana de 0,10 km entre a coordenada da usina e o nó
mais próximo. As oito covariáveis foram pré-declaradas antes do teste, e Holm
foi aplicado sobre as oito.

\begin{figure*}[t]
\centering
\includegraphics[width=\linewidth]{../figures/fig10_topologia_rede.pdf}
\caption{\textbf{A distância de rede até a demanda está associada à alocação do
corte que covariáveis locais não explicavam.}
\textbf{a}, Grafo de transmissão em operação, com subestações e linhas da EPE
separadas em 500 kV ou mais e abaixo disso, e as usinas fotovoltaicas da amostra
posicionadas no seu ponto de conexão. A cor codifica a taxa de corte na janela
de referência e o tamanho, a capacidade instalada. As estrelas marcam as nove
regiões metropolitanas usadas como centros de carga. A moldura cobre a região
que contém a amostra; 353 das 1.896 linhas ficam fora dela e não são desenhadas,
e todas as arestas participam do cálculo dos caminhos mínimos.
\textbf{b}, Correlação de Spearman entre a taxa de corte e cada uma das oito
covariáveis de rede pré-declaradas. Em vermelho as que sobrevivem à correção de
Holm.
\textbf{c}, Taxa de corte contra a distância de rede ao centro de carga mais
próximo, por subsistema, com a mediana por quartil de distância.
\textbf{d}, As duas covariáveis sobreviventes medidas dentro de cada subsistema,
para separar efeito de rede de efeito de subsistema.}
\label{fig:rede}
\end{figure*}

Duas covariáveis sobrevivem a Holm, com magnitude próxima e sinais opostos. A
distância de rede à carga dá $\rho = +0{,}38$ e a capacidade da usina dá
$\rho = -0{,}38$, ambas com $p = 0{,}023$ após correção. A soma de tensões no nó
atinge $p < 0{,}05$ bruto e não sobrevive à correção. As outras cinco não
atingem.

\conclusao{Quanto maior a distância em quilômetros de linha até o centro de
carga mais próximo, maior a taxa de corte. A estrutura local do nó não apresenta
associação.}

O resultado é consistente com a seção~\ref{sec:fig9}, em que o Nordeste é
cortado com o corredor operando abaixo do seu percentil habitual. As duas
medidas apontam para a distância até a demanda, e não para a saturação local do
ponto de conexão.

\campo{Não é apenas efeito de subsistema} A distância mediana à carga é próxima
nos dois subsistemas, com 497 km no Nordeste e 520 no Sudeste com Centro-Oeste.
A associação se mantém dentro do Nordeste, com $\rho = +0{,}34$, $p = 0{,}027$ e
$n = 43$. Ela não se mantém dentro do Sudeste com Centro-Oeste, onde dá
$\rho = +0{,}31$, $p = 0{,}25$ e $n = 16$. A capacidade da usina se mantém nos
dois, com $-0{,}41$ e $-0{,}44$, e no Nordeste com $p = 0{,}006$.

\campo{Uma sobrevivente sem hipótese direcional} A capacidade da usina tem sinal
negativo: as maiores são cortadas menos. A direção não havia sido prevista. As
explicações plausíveis, como qualidade de conexão e seleção contratual, são
testáveis e não foram testadas.

\begin{objecao}
A distância de rede não é elétrica: é caminho ponderado por quilômetro de linha.
A impedância está em dado aberto, em 100\% das linhas do pacote
\texttt{linha-transmissao}, com capacidade em MVA em 86,7\% delas. A escolha da
camada geométrica da EPE impediu o uso da distância elétrica nesta seção. Os
centros de carga são nove metrópoles, um proxy geográfico declarado. São oito
hipóteses sobre $n = 59$: Holm foi aplicado, e não há poder para efeitos
pequenos. O resultado é associação, e não causalidade.
\end{objecao}

\campo{Posição na literatura} Maji et al. \cite{maji2025} sustentam que o corte é
fenômeno de nó, e que estimá-lo exige descer ao ponto de conexão. É a hipótese
que esta seção testa e, na geração fotovoltaica, não confirma. O preditor é a
distância até a carga, uma grandeza de caminho, e não de nó. Newbery
\cite{newbery2024} fornece o enquadramento econômico: se a posição na rede
determina a exposição ao corte e o regime de acesso não a precifica, o sinal de
investimento aponta para o lugar errado.

\section{Replicação na geração eólica}
\label{sec:fig11}

\begin{ficha}
  \item[Janela]      novembro de 2022 a julho de 2025, eólica;
                     setembro de 2025 a agosto de 2026, solar
  \item[Amostra]     120 conjuntos eólicos; 59 solares
  \item[Unidade]     conjunto
  \item[Fonte]       ONS, \textit{constrained-off} eólico e fotovoltaico; EPE
  \item[Teste]       as mesmas 8 covariáveis, com Holm por fonte
\end{ficha}

O ONS apura \textit{constrained-off} eólico desde outubro de 2021, em 199
conjuntos, contra 87 do lado solar. A replicação usa as mesmas oito covariáveis,
na mesma ordem, sobre o mesmo grafo, com Holm reaplicado dentro de cada fonte.

\begin{figure*}[t]
\centering
\includegraphics[width=\linewidth]{../figures/fig11_replicacao_eolica.pdf}
\caption{\textbf{O resultado da seção anterior não se replica em eólica, e a
dissociação acompanha a origem declarada do corte de cada fonte.}
\textbf{a}, Correlação de Spearman entre a taxa de corte e as oito covariáveis
de rede pré-declaradas, medidas separadamente em fotovoltaica e eólica. A cor
cheia marca as que sobrevivem à correção de Holm aplicada dentro de cada fonte.
\textbf{b}, Energia cortada decomposta por razão e origem da restrição, conforme
classificação operativa do ONS, para cada fonte.
\textbf{c}, As duas covariáveis-chave medidas apenas no Nordeste, onde as duas
amostras coexistem, para separar tecnologia de geografia.
\textbf{d}, Taxa de corte eólica contra o número de conjuntos no mesmo ponto de
conexão, com a mediana por valor.}
\label{fig:eolica}
\end{figure*}

A distância de rede à carga, que na solar dá $+0{,}38$ e sobrevive a Holm, dá
$-0{,}07$ em 120 conjuntos eólicos. A amostra eólica é o dobro da solar, de modo
que a ausência não é falta de poder.

As covariáveis que sobrevivem em eólica são as que não sobreviveram em solar. O
número de usinas no mesmo nó vai de $+0{,}07$ para $+0{,}47$, e o grau do nó vai
de $-0{,}17$ para $+0{,}35$. As duas sobrevivem a Holm.

A classificação de origem do ONS, independente de todas as covariáveis testadas,
registra corte solar 91\% sistêmico e corte eólico 76\%, com 24\% de origem
local, ou 2,7 vezes a fração do solar. A confiabilidade local responde por 7\%
do corte solar e por 22\% do eólico.

\conclusao{A covariável que prevê o corte de cada fonte acompanha a origem
declarada do corte dessa fonte: grandeza sistêmica onde o corte é sistêmico,
grandeza local onde há corte local.}

\campo{Não é geografia} A amostra eólica é 91\% nordestina, e a solar não é.
Restringindo as duas ao Nordeste, a solar mantém distância à carga em $+0{,}34$
e usinas no nó em $+0{,}05$, enquanto a eólica apresenta $-0{,}19$ e $+0{,}36$.
A dissociação se mantém sob controle de geografia.

\campo{O que muda na seção anterior} A seção~\ref{sec:fig10} enunciou a hipótese
como sobre posição na rede, independentemente da tecnologia. A replicação
restringe o enunciado a fonte cujo corte é majoritariamente sistêmico. O
resultado permanece na geração fotovoltaica, e o que muda é a generalidade.

\campo{Consequência operativa} A exposição solar responde a aproximação da
demanda, por armazenamento ou deslocamento de carga. A exposição eólica responde
a reforço local de escoamento. Uma política única erraria em uma das duas.

\begin{objecao}
As janelas não coincidem: eólica de novembro de 2022 a julho de 2025, solar de
setembro de 2025 a agosto de 2026. Parte da diferença pode ser de período, e o
teste restrito ao Nordeste controla geografia, e não tempo. A eólica se aglomera
mais, com mediana de 4 conjuntos por nó contra 2 no solar, o que facilita
detectar aglomeração. A distância à carga, porém, tem faixa semelhante nas duas
fontes, e não prevê o corte eólico. Os regimes diferem nesta amostra, com 8\% de mediana
em eólica e 26\% em solar, contra 15,9\% e 22,1\% nas frotas completas. Origem é
classificação do ONS, e não inferência causal.
\end{objecao}

\campo{Posição na literatura} Bird, Lew e Milligan \cite{bird2016} organizam a
revisão de onze países pelo eixo entre excesso de oferta sistêmico e restrição
local. A EirGrid \cite{eirgrid} formaliza a distinção no vocabulário, com
\textit{curtailment} para razão sistêmica e \textit{constraint} para local, e
reporta 6,6\% contra 4,7\% no \textit{dispatch-down} eólico de 2025. Não foi
localizado teste publicado de se a covariável preditiva do corte varia com a
composição da causa, e a busca não foi sistemática.

\section{Robustez dos quatro resultados}
\label{sec:fig12}

\begin{ficha}
  \item[Janela]      as amostras das duas seções anteriores
  \item[Amostra]     59 conjuntos solares; 120 eólicos
  \item[Unidade]     conjunto
  \item[Fonte]       ONS, Rede de Operação por número de barra; EPE
  \item[Teste]       bootstrap com 4.000 reamostras; $z$ de Fisher
\end{ficha}

O ONS publica a rede com conectividade por número de barra, impedância em 100\%
das linhas e capacidade em MVA em 87\%. Isso permite reconstruir o grafo por via
independente da camada geométrica usada nas duas seções anteriores. Os três
critérios são o intervalo bootstrap que exclui zero, a estabilidade entre as
metades de uma partição estrutural e a reprodução no grafo alternativo.

\begin{figure*}[t]
\centering
\includegraphics[width=\linewidth]{../figures/fig12_robustez.pdf}
\caption{\textbf{Dos quatro resultados de rede, um passa nos três critérios e o
principal da seção~\ref{sec:fig10} passa em um.}
\textbf{a}, Distribuição bootstrap do coeficiente de Spearman de cada resultado,
com 4.000 reamostras com reposição. A barra marca o intervalo de 95\% e o ponto,
o valor da amostra.
\textbf{b}, Coeficiente calculado separadamente nas duas metades de uma partição
estrutural, entre usinas cuja barra de conexão consta do cadastro de rede de
230 kV ou mais do ONS e as demais, com o valor de $p$ do teste $z$ de Fisher
para a diferença entre metades.
\textbf{c}, O mesmo coeficiente calculado no grafo da EPE e no grafo do ONS
construído por número de barra. Abaixo de cada barra, o tamanho da amostra no
grafo alternativo, ou a marca de que o critério não se aplica.
\textbf{d}, Quantos dos critérios aplicáveis cada resultado passa. O traço marca
o critério não aplicável.}
\label{fig:robustez}
\end{figure*}

\begin{table}[t]
\centering
\tabfont
\begin{tabular}{@{}>{\rag\arraybackslash}p{3.4cm}ccc@{}}
\toprule
\cab{Resultado} & \cab{IC} & \cab{estável} & \cab{outro grafo} \\
\midrule
eólica, grau do nó         & sim & sim & sim \\
eólica, usinas no nó       & sim & não & sim \\
solar, capacidade da usina & sim & sim & --- \\
solar, distância à carga   & sim & não & não \\
\bottomrule
\end{tabular}
\caption{Critérios de robustez aplicados aos quatro resultados que sobreviveram
a Holm.}
\label{tab:robustez}
\notas{O traço indica critério não aplicável: a capacidade da usina não deriva
do grafo e não tem versão alternativa. Para ela valem dois critérios, e não
três.}
\end{table}

Sobreviveram a Holm quatro resultados, e não três. A capacidade da usina em
solar, com $\rho = -0{,}38$ e $p = 0{,}023$ após correção, tem magnitude igual à
da distância à carga e não constava do conjunto avaliado antes.

A distância à carga tem intervalo bootstrap de 0,12 a 0,61, uma largura de 0,49,
cerca do dobro dos demais. Partindo a amostra, o coeficiente cai de 0,62 para
0,11 conforme a metade, com $p = 0{,}025$, e a metade de maior coeficiente
contém todas as 16 usinas do Sudeste. Isso reabre o confundimento regional que o
painel \textbf{d} da seção~\ref{sec:fig10} afastava.

A capacidade da usina passa nos dois critérios aplicáveis, com intervalo de
$-0{,}60$ a $-0{,}13$ e partição não rejeitada, de $-0{,}18$ contra $-0{,}48$ e
$p = 0{,}21$.

\conclusao{O grau do nó na geração eólica passa nos três critérios. A distância
à carga na geração fotovoltaica passa em um dos três, e é o resultado de rede
menos sustentado.}

Dois critérios não são três, e a diferença entre metades ainda reduz o
coeficiente da capacidade da usina em dois terços. Passar nos critérios
aplicáveis é não ser rejeitado, e não é ser confirmado.

\campo{Amostra, e não método} No grafo do ONS a distância dá $-0{,}14$ em
quilômetros e $-0{,}06$ em reatância. Rodando o grafo da EPE nas mesmas 28
usinas, ele também não apresenta associação. A diferença é de amostra, o que não
recupera o resultado e identifica em que subconjunto ele ocorre.

A distância elétrica foi calculada, e a afirmação da seção~\ref{sec:fig10} de
que ela não estaria disponível é incorreta. A correlação entre a distância
elétrica e a distância em quilômetros é de 0,53, de modo que a substituição
constitui teste independente. O resultado não muda.

\campo{Os critérios não dependiam de fonte nova} Os três critérios usam apenas
dados já disponíveis na seção~\ref{sec:fig10}, e poderiam ter sido aplicados
ali. Foram executados a partir da troca de grafo.

\begin{objecao}
A partição não é aleatória, por construção: ela testa dependência a
característica estrutural, e não variabilidade amostral genérica. O grafo do ONS
é exato onde cobre, e exclui conexões abaixo de 230 kV; para a amostra solar
isso remove mais da metade dos conjuntos e elimina um subsistema inteiro. Três
critérios não esgotam robustez, e faltam validação temporal, análise de
influência ponto a ponto e sensibilidade à definição de centro de carga. Os
limiares são convenção declarada, e o critério do intervalo foi corrigido de
``é positivo'' para ``exclui zero''. Como o grafo é variável derivada de
sensoriamento, ele propaga erro para o intervalo \cite{proctor2023}, e nenhuma
imputação foi aplicada.
\end{objecao}

\chapter{Produção e exposição ao corte de uma área nova}
\label{cap:aplicacao}

Estimar quanto uma área nova junto a Sol do Cerrado produziria, e quanto dessa
energia o despacho absorveria, depende de três fatores. Os três já foram
medidos nos capítulos anteriores, e nenhum deles é suposto.

\section{A cadeia de conversão}

\begin{table}[t]
\centering
\tabfont
\begin{tabular}{@{}>{\rag\arraybackslash}p{3.3cm}
                  >{\rag\arraybackslash}p{2.1cm}
                  >{\rag\arraybackslash}p{1.9cm}@{}}
\toprule
\cab{Passo} & \cab{Fator} & \cab{Origem} \\
\midrule
terreno para módulo     & dividir por 2,5
                        & seção~\ref{sec:fig3} \\
módulo para capacidade  & 5,12 m\textsuperscript{2}/kWp
                        & seção~\ref{sec:fig2} \\
capacidade para potencial & fator de 0,229
                        & seção~\ref{sec:fig6} \\
\bottomrule
\end{tabular}
\caption{Fatores da cadeia e a análise que produz cada um.}
\label{tab:cadeia}
\notas{O fator de capacidade potencial vem da razão de desempenho limpa de Sol
do Cerrado, de 0,931, calibrada em 279 dias sem restrição, e do valor de
irradiância da NASA POWER sobre Jaíba, de 5,908 kWh/m\textsuperscript{2} por
dia.}
\end{table}

Aplicada à usina existente, a cadeia fecha. Os 681,3 MW fiscalizados implicam
3,49 km\textsuperscript{2} de módulo, e o caminho inverso devolve 681,3 MW. Isso
não valida a cadeia. É divisão e multiplicação pela mesma constante de
5,12 m\textsuperscript{2}/kWp, e fecharia com qualquer valor dela. A ida e volta
confere a aritmética.

O elo não validado é o primeiro. A razão entre área de terreno e área de módulo,
de 2,5, é típica de projeto e não foi medida nesta usina. Ela entra linearmente
na capacidade, de modo que um erro de 20\% nela produz 20\% de erro na
capacidade. Medi-la exigiria delinear o arranjo por segmentação, o passo que a
seção~\ref{sec:fig7} deixa em aberto.

\section{O que a cadeia estima, e contra o que ela é conferida}

Para 2025 o método prevê fator de capacidade de 0,229 e 1.368 GWh. A usina
entregou fator de 0,150 e 897 GWh. A diferença é de 34,4\%.

A apuração do ONS dá 32,8\% para a mesma usina e o mesmo ano, por outro
caminho: curva de desempenho declarada da usina, e apuração operativa de
\textit{constrained-\allowbreak off}. As duas diferem em 1,6 ponto.

\conclusao{Uma estimativa de perda construída a partir de irradiância e razão de
desempenho calibrada chega a 1,6 ponto da apuração operativa do operador.}

\campo{A independência é parcial} A cadeia não usa a geração de referência do
operador, e essa é a parte externa. Ela usa a razão de desempenho limpa de
0,931, calibrada nos dias que a própria apuração do ONS marca como sem
restrição. O sinalizador de restrição é compartilhado. O que se testa é a
estimativa de quanto, e não a de quando.

Não é uma terceira estimativa. A razão $1 - PR_{2025}/PR_{\text{limpo}}$ é o
mesmo estimador da seção~\ref{sec:fig6}, que dá 33,5\%, e a diferença para os
34,4\% vem da agregação da irradiância. São duas estimativas, e o que este
capítulo acrescenta é o passo de área para capacidade, que é o elo não validado.

\section{Cenários para uma área nova}

\begin{table}[t]
\centering
\tabfont
\begin{tabular}{@{}rrrrr@{}}
\toprule
\cab{AOI} & \cab{capac.} & \cab{potencial} & \cab{a 31\%} & \cab{a 62\%} \\
km\textsuperscript{2} & MW & GWh/ano & GWh & GWh \\
\midrule
 1 &    78 &   157 &   108 &    60 \\
 2 &   156 &   314 &   216 &   119 \\
 5 &   391 &   784 &   541 &   298 \\
10 &   781 & 1.569 & 1.082 &   596 \\
20 & 1.563 & 3.137 & 2.165 & 1.192 \\
\bottomrule
\end{tabular}
\caption{Produção estimada de uma área nova em Jaíba, sob duas taxas de corte.}
\label{tab:cenarios}
\notas{A coluna de 31\% aplica a taxa média medida no ponto. A coluna de 62\%
aplica a razão entre taxa marginal e taxa média medida por Novan e Wang
\cite{novan2024}.}
\end{table}

As duas últimas colunas separam taxa média de taxa marginal. A taxa média não se
aplica a uma usina nova. Novan e Wang \cite{novan2024} mediram que a taxa
marginal, que é quanto de uma usina nova é cortado, fica em cerca de duas vezes
a média.

\section{Exposição ao corte no ponto de conexão}

O barramento de Jaíba tem 1.301 MW em três conjuntos: Sol do Cerrado com 681 MW,
Jaíba V com 500 MW e Jaíba 138 kV com 120 MW. Todos são cortados entre 24\% e
31\% na janela de referência.

O corte no ponto é 89,5\% de origem sistêmica, com 61,8\% por razão energética,
16,2\% por confiabilidade e 11,5\% por indisponibilidade externa. Apenas 10,5\%
tem origem local. No agregado nacional a distribuição é de 68,5\% energética,
22,0\% confiabilidade e 9,5\% indisponibilidade externa. A restrição não é de
saturação do barramento, e sim de oferta excedente no instante do corte.

O resultado é consistente com os capítulos~\ref{cap:despacho}
e~\ref{cap:rede}: térmica inflexível nas horas de corte, corredor abaixo do seu
percentil habitual e distância à demanda como preditor. Jaíba está a 605 km de
rede do centro de carga mais próximo, no terceiro quartil da amostra.

\campo{Uma leitura que os dados não sustentam} A hipótese de que mais capacidade
no mesmo nó produz mais corte não se sustenta. Na seção~\ref{sec:fig10} a
capacidade renovável no nó dá $\rho = -0{,}24$ e não sobrevive a Holm, e o mesmo
vale para o grau, a tensão e o número de usinas. O que sobrevive é a distância.
A leitura defensável é a de posição na rede, e não a de saturação do ponto. Uma
usina nova ali herdaria essa posição.

\begin{objecao}
O cálculo não usa limite de linha. O ONS publica capacidade operativa em MVA,
com variantes sazonais, em 86,7\% das linhas, e incorporá-la tornaria a
estimativa de exposição mais precisa. A razão de 2 vezes entre taxa marginal e
média é medida na Califórnia; não há equivalente brasileiro publicado, e ela
entra como ordem de grandeza declarada, e não como fator calibrado. A razão
entre terreno e módulo de 2,5 é o elo não validado da cadeia, e inclinação,
rastreamento e relevo a deslocam. O fator de capacidade potencial usa o
desempenho limpo de Sol do Cerrado; uma usina nova com outro módulo ou outro
rastreamento teria outro.
\end{objecao}

\campo{Escopo} O resultado é uma estimativa física e operativa: quanto a área
produziria e quanto o despacho absorveria. Ela cobre dimensionamento, comparação
de sítios e antecipação de exposição ao corte. Contrato, outorga e parecer de
acesso não estão em dado aberto e não entram.

{\color{tinta2}\footnotesize Dados-fonte em
\texttt{data/processed/aplicacao\_*.csv}, gerados por
\texttt{notebooks/\_aplicacao\_expansao.py}.}


\part{Síntese}
\chapter{O que o conjunto sustenta}
\label{cap:estabelece}

Nove resultados se sustentam, oito conclusões foram retiradas por análise
posterior e três perguntas ficam em aberto. Cada uma das conclusões retiradas
foi afirmada, em algum momento, com a mesma confiança das que ficaram.

\section{Magnitude e composição do corte}

O corte é grande e de origem sistêmica. A taxa é de 21,6\% no parque com
apuração, dois terços por razão energética, incidindo sobre toda a janela solar.

A geração de referência do ONS é calibrada. O viés mediano é de 1,02 em 74
usinas, e não se degrada onde o corte é mais intenso. O agregado da
seção~\ref{sec:fig1} permanece válido.

A perda de Sol do Cerrado é de rede, e não de recurso. Nove décimos da queda de
29\% no fator de capacidade não se explicam pela irradiância.

O corte pode ser estimado sem a referência do operador. Irradiância e razão de
desempenho calibrada estimam 33,5\% de perda em 2025, e 34,4\% com o passo de
área acoplado, contra 32,8\% apurados pelo ONS.

\begin{objecao}
A independência dessa estimativa é parcial. Ela dispensa a geração de
referência, e ainda usa a apuração do ONS para selecionar os dias de calibração.
O sinalizador de restrição é compartilhado entre as duas.
\end{objecao}

\section{Observabilidade por satélite}

A geração distribuída é sub-pixel. São 91\% das instalações abaixo de um pixel
Sentinel-2, e a verificação por imagem exige 2 m ou melhor, além de coordenada
por instalação, que não existe no registro.

A construção de usina centralizada é detectável e datável; a operação não é. O
brilho relativo sobe 34\% na construção, com $p = 1{,}2 \times 10^{-5}$, e não
apresenta sinal na operação, com $p = 0{,}09$.

\section{Despacho e alocação}

O piso térmico inflexível está no Sudeste, e não no Nordeste. São 1.739 MWmed
de inflexibilidade pura no primeiro, contra 36 no segundo.

O grau do nó está associado ao corte eólico, com $\rho = +0{,}35$. É o único dos
quatro resultados de rede a passar em bootstrap, partição estrutural e troca de
grafo.

O preditor acompanha a origem do corte. A geração solar tem corte 91\% sistêmico
e responde à distância até a demanda. A eólica tem 24\% de corte local e responde
a aglomeração no nó, com $\rho = +0{,}47$. A dissociação se mantém sob controle
de geografia.

\section{As conclusões retiradas}

As oito afirmações abaixo foram sustentadas por uma análise e retiradas por
outra. Elas permanecem no relatório porque a correção é parte do resultado.

\begin{enumerate}[leftmargin=1.2em,itemsep=0.3em]
\item \textbf{A referência do ONS superestima o corte.} Era artefato do
contrafactual, porque 2023 já era um ano com restrição. A revisão reforça a
conclusão principal: o corte implícito é de 33\%, e não de 26\%.

\item \textbf{A térmica inflexível desloca a geração solar.} Essa térmica se
associa igualmente ao corte a dois mil quilômetros de distância. Ela é marcador
sistêmico, e o desenho não separa as hipóteses.

\item \textbf{O Nordeste é cortado porque o corredor satura.} Nas horas de corte
alto o fluxo para o Sudeste é menor que nas demais.

\item \textbf{A térmica não abre espaço para a solar.} A afirmação não sobrevive
ao controle de carga, porque a carga líquida diária é plana.

\item \textbf{Explicações locais do corte generalizam.} Nenhuma das cinco
covariáveis testadas no par se reproduz no parque.

\item \textbf{A distância à demanda vale para qualquer tecnologia.} A replicação
em eólica restringe o enunciado a fonte de corte sistêmico, e os critérios de
robustez mostram que, mesmo na solar, o resultado é o menos sustentado dos
quatro.

\item \textbf{Impedância e limite de linha não estão em dado aberto.} Estão, no
pacote \texttt{linha-transmissao} do ONS, com reatância em 100\% das linhas e
capacidade em MVA em 87\%. A afirmação incorreta aparecia em quatro pontos.

\item \textbf{Detectar painéis por satélite audita política de geração
distribuída.} Era a premissa original, e a seção~\ref{sec:fig2} a encerra.
\end{enumerate}

\chapter{O que falta}
\label{cap:falta}

\section{Perguntas em aberto}

\campo{Alocação do corte entre usinas} A seção~\ref{sec:fig10} identificou duas
covariáveis que sobrevivem à correção de multiplicidade: distância de rede à
demanda e capacidade da usina. Sob os critérios da seção~\ref{sec:fig12} é a
segunda que resiste, e é a que não tem explicação testada. A variância explicada
é pequena, e a maior parte da dispersão da taxa de corte permanece sem
explicação.

\campo{Causa do corte no Nordeste} Nem piso térmico local, nem escoamento. Resta
ausência de demanda no destino, obtida por eliminação e não testada de forma
positiva.

\campo{Inflexibilidade declarada contra mínimo técnico real} É a grandeza que
uma auditoria exigiria, e não está publicada.

\section{Trabalho seguinte, em ordem de retorno esperado}

\begin{enumerate}[leftmargin=1.2em,itemsep=0.35em]
\item \textbf{Reconstruir o corte anterior a abril de 2024.} A maquinaria da
seção~\ref{sec:fig6} já existe. Um modelo treinado nos intervalos sem restrição,
aplicado de forma retroativa, viabiliza o teste temporal da
seção~\ref{sec:fig4}. Há verdade de campo para validar, nos dois anos em que as
duas séries coexistem.

\item \textbf{Medir o resíduo não explicado.} No nível de intervalo, a diferença
de desempenho entre um modelo só com observáveis sistêmicos e outro com
identidade de usina quantifica o que os capítulos de rede deixaram sem
explicação. A série eólica acrescenta 1,4 GB e três anos de painel. A primeira
metade já está publicada: Gorka e Roald \cite{gorka2024} classificam a
ocorrência de corte no CAISO a partir do estado do sistema, que é o modelo sem
identidade. O que não se encontrou publicado é o passo seguinte, de recalcular o
desempenho \emph{dentro} de cada intervalo, onde a informação de quando o corte
ocorre é removida por construção e resta a de qual usina o recebe.

\item \textbf{Substituir NASA POWER por NSRDB.} A célula passa de cerca de 50 km
para 2 km, e o passo de diário para 10 minutos, sem custo. Atende três das
objeções já escritas.

\item \textbf{Delinear o arranjo por segmentação.} Substitui os círculos de
400 m da seção~\ref{sec:fig7}, aumenta a precisão da datação e mede a razão
entre terreno e módulo que o capítulo~\ref{cap:aplicacao} assume.
\end{enumerate}

\section{Uma regra de método}

Qualquer saída de modelo que entre em inferência a jusante enviesa o ponto e o
erro padrão \cite{proctor2023}. A predição entra como distribuição, e não como
valor pontual. A regra vale para a irradiância do capítulo~\ref{cap:calibracao}
e para o grafo do capítulo~\ref{cap:rede}, e não foi aplicada em nenhum dos
dois.

\campo{O cálculo está disponível} Angelopoulos et al. \cite{angelopoulos2023}
dão o procedimento: intervalo de confiança válido para média, quantil e
coeficiente de regressão, sem hipótese sobre o algoritmo que produz a predição,
desde que exista um subconjunto com o valor observado. Kluger et al.
\cite{kluger2025} estendem o procedimento ao caso em que a predição entra como
covariável imputada, e motivam com predição de indicador ambiental por imagem de
satélite alimentando regressão a jusante, que é a forma exata dos capítulos de
calibração e de rede.

A condição é ter verdade de campo num subconjunto. Para a irradiância ela
existe: a série medida no plano do arranjo cobre a mesma janela. Para o grafo
não há equivalente publicado, e ali a correção exige antes definir contra o quê
a predição seria comparada.

\begin{objecao}
Aplicá-la alargaria os intervalos de todos os resultados de rede. Nenhum deles
foi medido sob essa correção, e não se pode afirmar que o resultado que passa
nos três critérios sobreviveria a ela.
\end{objecao}


\appendix
\onecolumn
% Apendice em pagina deitada. A `landscape` abre pagina, entao ela
% comeca ANTES do titulo: assim o titulo encabeca a tabela em vez de
% ficar orfao numa pagina em pe. Fora do regime de duas colunas, porque
% `longtable` exige uma coluna so.
\clearpage
\begin{landscape}

\chapter{Inventário das análises}
\label{ap:inventario}

\begin{promancha}
Cada linha declara a janela, o tamanho da amostra, a unidade que se repete
dentro dela e o teste aplicado. As janelas de corte saem da mesma regra,
percentil 80 do início e percentil 20 do fim da apuração, e por isso diferem
entre fontes de geração. A análise 1 vem de extração anterior às demais.
\end{promancha}

\tabfont
\begin{longtable}{@{}
    >{\rag\arraybackslash}p{0.7cm}
    >{\rag\arraybackslash}p{3.9cm}
    >{\rag\arraybackslash}p{3.7cm}
    >{\rag\arraybackslash}p{3.0cm}
    >{\rag\arraybackslash}p{3.1cm}
    >{\rag\arraybackslash}p{5.5cm}
    >{\rag\arraybackslash}p{1.3cm}@{}}
\toprule
\cab{\#} & \cab{Objeto} & \cab{Janela} & \cab{Amostra} &
\cab{Unidade de repetição} & \cab{Teste declarado} & \cab{Seção} \\
\midrule
\endfirsthead
\toprule
\cab{\#} & \cab{Objeto} & \cab{Janela} & \cab{Amostra} &
\cab{Unidade de repetição} & \cab{Teste declarado} & \cab{Seção} \\
\midrule
\endhead
\bottomrule
\endfoot
1 & Corte fotovoltaico no SIN & abr/2024 a jul/2026 & 85 usinas & usina por 30 min & nenhum: descritiva & \ref{sec:fig1} \\
2 & Detectabilidade da geração distribuída & conexões até abr/2026 & 4,6 milhões de instalações & empreendimento & nenhum: descritiva & \ref{sec:fig2} \\
\addlinespace[.5ex]
3 & Estudo de caso pareado & out/2024 a ago/2026 & 2 ativos; 60 conjuntos & conjunto & Spearman sobre 5 covariáveis, sem correção & \ref{sec:fig3} \\
4 & Cronologia do entorno & 2023 a jul/2026; entorno desde 2021 & 43 e 24 meses & conjunto por mês & nenhum teste formal & \ref{sec:fig4} \\
\addlinespace[.5ex]
5 & Controle climático & 2023 a jul/2026 & meses com cobertura climática & conjunto por mês & nenhum: decomposição aritmética & \ref{sec:fig5} \\
6 & Calibração da referência & abr/2024 a jul/2026 & 74 usinas calibradas & usina; dia-usina & Spearman sobre 2 covariáveis & \ref{sec:fig6} \\
\addlinespace[.5ex]
7 & Detecção de construção & 2018 a 2026 & 34 conjuntos, 306 medições & conjunto por ano & Wilcoxon pareado, 2 comparações & \ref{sec:fig7} \\
\addlinespace[.5ex]
8 & Despacho fora da ordem de mérito & abr/2024 a ago/2026 & 21.192 horas & hora do SIN & Mann-Whitney e Spearman, por quintil & \ref{sec:fig8} \\
9 & Decomposição por subsistema & idem, 9 h a 15 h & 21.192 horas & hora do subsistema & Spearman, sem correção & \ref{sec:fig9} \\
\addlinespace[.5ex]
10 & Topologia da rede, solar & set/2025 a ago/2026 & 59 conjuntos & conjunto & Spearman sobre 8 covariáveis, com Holm & \ref{sec:fig10} \\
11 & Replicação em eólica & nov/2022 a jul/2025 eólica; set/2025 a ago/2026 solar & 120 e 59 conjuntos & conjunto & idem, com Holm por fonte & \ref{sec:fig11} \\
12 & Robustez dos resultados de rede & amostras de 10 e 11 & 59 e 120 conjuntos & conjunto & bootstrap com 4.000 reamostras; $z$ de Fisher & \ref{sec:fig12} \\
\addlinespace[.5ex]
aux & Imagem comercial & catálogo de 2026 & 21 sensores & sensor & nenhum: apoio de decisão & \ref{sec:aux} \\
apl & Área nova em Jaíba & 2025 & 1 ponto de conexão & cenário de área & nenhum: propagação de fatores & --- \\
\end{longtable}

\begin{promancha}
\vspace{0.6em}
Somados os testes que vão a arquivo, o relatório reporta 49: cinco na análise 3,
dois na 7, cinco na 8, catorze na 9, oito na 10, oito na 11 e sete na 12. Desses,
16 sobrevivem à correção de Holm aplicada dentro da análise que os produz, e 33
são reportados sem correção. Não há correção entre análises, e um limiar em
nível de relatório não é aplicado.
\end{promancha}

\end{landscape}

\clearpage
\begin{landscape}

\chapter{Procedência das fontes}
\label{ap:procedencia}

\begin{promancha}
Todas as análises usam dados abertos e sem autenticação. As fontes não têm o
mesmo peso probatório, e a tabela final separa as revisadas por pares das
demais.
\end{promancha}

\section*{Reprodução}

\begin{promancha}
Cada figura exporta em CSV os dados que produzem cada painel, sob
\texttt{data/processed/}. Os notebooks e scripts geradores ficam em
\texttt{notebooks/}, e as figuras compostas em \texttt{figures/}. As métricas
citadas no corpo saem desses arquivos, e não da leitura da imagem.
\end{promancha}

\section*{Bases de dados}

\begin{promancha}
\campo{ONS, Portal de Dados Abertos} Restrição por \textit{constrained-off}
fotovoltaico e eólico; geração térmica por motivo de despacho; fator de
capacidade; curva de carga; intercâmbio nacional; energia armazenada por
subsistema; Custo Marginal de Operação semi-horário; relação entre conjunto e
usina; linhas de transmissão e subestações da Rede de Operação.

\campo{ANEEL} Sistema de Informações de Geração; relação de empreendimentos de
geração distribuída e informações técnicas fotovoltaicas; SIGET, com contrato,
módulo, linha de transmissão e subestações de origem e destino.

\campo{EPE} Webmap institucional, em serviço ArcGIS REST: camadas de linhas de
transmissão e subestações em operação, com geometria.

\campo{ESA e Copernicus} Sentinel-2 L2A, pelo catálogo STAC público da AWS em
\texttt{earth-\allowbreak search.\allowbreak aws.\allowbreak element84.\allowbreak com}.

\campo{NASA POWER} Irradiância diária, de MERRA-2 e SYN1DEG.

\campo{IBGE} Malhas territoriais das unidades da federação.
\end{promancha}

\section*{Comparação internacional}

\begin{promancha}
\tabfontw
\begin{tabular}{@{}>{\rag\arraybackslash}p{6.4cm}rr
                  >{\rag\arraybackslash}p{5.0cm}@{}}
\toprule
\cab{Sistema} & \cab{Taxa} & & \cab{Referência} \\
\midrule
Brasil, solar com apuração    & 21,6\% & & este relatório \\
Brasil, eólica e solar, 2025  & 20,7\% & & ONS \\
Chile, 2024                   & 17\%   & & imprensa setorial \\
Brasil, eólica com apuração   & 15,9\% & & este relatório \\
Irlanda, eólica, 2025         & 11,3\% & & EirGrid \\
Espanha, julho de 2025        & 11\%   & & imprensa setorial \\
Califórnia, solar             &  4,3\% & & Novan e Wang \\
\bottomrule
\end{tabular}
\notas{Os denominadores diferem, e a comparação com sistemas estrangeiros é de
ordem de grandeza, e não pareada. As duas linhas brasileiras são pareadas entre
si: ambas são razão de energia cortada sobre disponível, na janela de abril de
2024 a agosto de 2026, em que solar e eólica têm apuração simultânea, sobre 87
conjuntos solares e 178 eólicos. Numa extração única elas dão 22,1\% e 15,9\%; a
linha da solar preserva os 21,6\% da seção~\ref{sec:fig1}, de extração
anterior.}
\end{promancha}

\begin{promancha}
As taxas brasileiras ficam no topo da faixa internacional para as duas fontes. A
solar perde cerca de 1,4 vez o que a eólica perde. O que separa os dois regimes
não é a magnitude do corte, e sim a composição da causa: 91\% de origem
sistêmica na solar contra 76\% na eólica, medido na seção~\ref{sec:fig11}.
\end{promancha}

\section*{Mudança regulatória durante a execução}

\begin{promancha}
A pergunta que a seção~\ref{sec:fig3} deixou em aberto, sobre se o corte é
ressarcido, deixou de ser hipotética. A Lei 15.269/2025 criou compensação por
\textit{constrained-off} para eólicas e solares, custeada por Encargos de
Serviços do Sistema, com efeitos retroativos a 1 de setembro de 2023 e
condicionada à renúncia de ações judiciais e à assinatura de termo de
compromisso. A Portaria MME 140 regulamentou os cortes ocorridos até 25 de
novembro de 2025, e os posteriores ficam com a ANEEL.

A primeira fase reuniu manifestação de interesse de 1.539 usinas, com cerca de
53 GW instalados. Estimativas setoriais apontam perda da ordem de R\$ 6,5
bilhões em 2025, mesma ordem de grandeza da valoração ao CMO da
seção~\ref{sec:fig3} extrapolada para o parque.

Isso reforça a objeção daquela seção. O CMO mede valor sistêmico, e não
ressarcimento. Com a lei, os dois divergem por construção: o gerador é
compensado por energia cujo valor sistêmico, no instante do corte, era próximo
de zero.
\end{promancha}

\section*{Peso das fontes}

\begin{promancha}
\campo{Revisadas por pares} Vieira et al. (2026), Novan e Wang (2024),
Kruitwagen et al. (2021), Yu et al. (2018), Jiang et al. (2021), Alix-García e
Millimet (2023), Bird et al. (2016), Callaway e Sant'Anna (2021), Maji et al.
(2025).

\campo{Working papers e pré-prints} Newbery (2024), Proctor et al. (2023), Silva
et al. (2026), Global Renewables Watch (2025), Robinson et al. (2021).

\campo{Fontes normativas} Lei 15.269/2025, Portaria MME 140 e Lei 14.300/2022.

\campo{Imprensa setorial} Taxas de Chile, Espanha e Europa, e o valor de R\$ 6,5
bilhões. Usadas como ordem de grandeza, e nunca como evidência de mecanismo.
\end{promancha}

\end{landscape}


\backmatter
\twocolumn
\begin{thebibliography}{99}
\footnotesize

\bibitem{alixgarcia2023} Alix-García, J. e Millimet, D.L. Remotely incorrect?
  Accounting for nonclassical measurement error in satellite data on
  deforestation. \emph{Journal of the Association of Environmental and Resource
  Economists}, 10(5), 1335--1367, 2023.

\bibitem{angelopoulos2023} Angelopoulos, A.N., Bates, S., Fannjiang, C.,
  Jordan, M.I. e Zrnic, T. Prediction-powered inference. \emph{Science},
  382(6671), 669--674, 2023.

\bibitem{bird2016} Bird, L., Lew, D., Milligan, M. et al. Wind and solar energy
  curtailment: a review of international experience. \emph{Renewable and
  Sustainable Energy Reviews}, 65, 577--586, 2016.

\bibitem{callaway2021} Callaway, B. e Sant'Anna, P.H.C. Difference-in-differences
  with multiple time periods. \emph{Journal of Econometrics}, 225(2), 2021.

\bibitem{cem2018} Clean Energy Ministerial. \emph{Thermal power plant
  flexibility}. Advanced Power Plant Flexibility Campaign, 2018.

\bibitem{eirgrid} EirGrid. \emph{Annual renewable energy constraint and
  curtailment report}, edições de 2013 a 2024.

\bibitem{gorka2024} Gorka, J. e Roald, L. Classification models for
  forecasting and real-time identification of solar curtailment in the
  California grid. \emph{ACM e-Energy '24}, Singapura, 2024.

\bibitem{grw2025} Microsoft Research. Global Renewables Watch: a temporal
  dataset of solar and wind energy derived from satellite imagery.
  arXiv:2503.14860, 2025.

\bibitem{jiang2021} Jiang, H. et al. Multi-resolution dataset for photovoltaic
  panel segmentation from satellite and aerial imagery. \emph{Earth System
  Science Data}, 13, 5389--5401, 2021.

\bibitem{kluger2025} Kluger, D.M., Lu, K., Zrnic, T., Wang, S. e Bates, S.
  \emph{Prediction-powered inference with imputed covariates and nonuniform
  sampling}. arXiv:2501.18577, 2025.

\bibitem{kruitwagen2021} Kruitwagen, L., Story, K.T., Friedrich, J., Byers, L.,
  Skillman, S. e Hepburn, C. A global inventory of photovoltaic solar energy
  generating units. \emph{Nature}, 598, 604--610, 2021.

\bibitem{maji2025} Maji, D., Irwin, D.E., Shenoy, P.J. e Sitaraman, R.K. A first
  look at node-level curtailment of renewable energy and its implications.
  \emph{ACM e-Energy '25}, Roterdã, 2025.

\bibitem{newbery2024} Newbery, D. \emph{Marginal curtailment of wind and solar
  PV: transmission constraints, pricing and access regimes for efficient
  investment}. EPRG Working Paper 2401, University of Cambridge, 2024.
  Ver também \emph{Marginal curtailment under network constraints}, Cambridge
  Working Paper in Economics 2581, 2025.

\bibitem{novan2024} Novan, K. e Wang, Y. Estimates of the marginal curtailment
  rates for solar and wind generation. \emph{Journal of Environmental Economics
  and Management}, 2024.

\bibitem{oshaughnessy2020} O'Shaughnessy, E., Cruce, J. e Xu, K. Too much of a
  good thing? Global trends in the curtailment of solar PV. \emph{Solar Energy},
  208, 1068--1077, 2020.

\bibitem{proctor2023} Proctor, J., Carleton, T. e Sum, S. \emph{Parameter
  recovery using remotely sensed variables}. NBER Working Paper 30861, 2023.

\bibitem{robinson2021} Robinson, C. et al. Temporal cluster matching for change
  detection of structures from satellite imagery. Microsoft Research, 2021.

\bibitem{silva2026} Nascimento, P., Cossich, W., Araujo, L., Santos, I.,
  Almeida, K. e Marcato, A. An integrated methodology for assessing wind power
  curtailment using anemometric measurements and operational data in the
  Brazilian context. \emph{Atmosphere}, 17(4), 333, 2026.

\bibitem{vieira2026} Vieira, G.T.T., Silva, W.N., Lourenço, L.F.N., Monaro,
  R.M., Salles, M.B.C. e Almeida, C.F.M. Characterizing wind and solar
  curtailment in Brazil: an evidence-based analysis of operational drivers.
  \emph{Electric Power Systems Research}, 2026.

\bibitem{yu2018} Yu, J., Wang, Z., Majumdar, A. e Rajagopal, R. DeepSolar: a
  machine learning framework to efficiently construct a solar deployment
  database in the United States. \emph{Joule}, 2(12), 2018.

\end{thebibliography}


\end{document}
